% part: intuitionistic-logic
% chapter: introduction
% section: syntax

\documentclass[../../../include/open-logic-section]{subfiles}

\begin{document}

\olfileid[pt]{int}{int}{syn}

\olsection{Sintaxe da lógica intuicionista}

A sintaxe da lógica intuicionista é a mesma da lógica proposicional.
Na lógica proposicional clássica, é possível definir conectivos por
meio de outros: p.ex., pode-se definir $!A \lif !B$ por
$\lnot !A \lor !B$, ou $!A \lor !B$ por $\lnot(\lnot !A \land \lnot
!B)$. Por isso, apresentações da lógica clássica frequentemente
introduzem alguns conectivos como abreviações dessas definições. Isso
não ocorre na lógica intuicionista, com duas exceções: $\lnot !A$ pode
ser---e muitas vezes é---definida como abreviação de $!A \lif \lfalse$.
Nesse caso, é claro, $\lfalse$ não pode ser ela própria definida!{}
Além disso, $!A \liff !B$ pode ser definida, como na lógica clássica,
por $(!A \lif !B) \land (!B \lif !A)$.

As !!^{formula}s da lógica proposicional intuicionista são construídas
a partir de \emph{!!{propositional variable}s} e da constante
proposicional~$\lfalse$, mediante \emph{conectivos lógicos}. Temos:

\begin{enumerate}
\item Um conjunto~$\PVar$ !!^a{denumerable} de !!{propositional variable}s $\Obj p_0$,
  $\Obj p_1$, \dots
  \item A constante proposicional de !!{falsity}~$\lfalse$.
\item Os conectivos lógicos: $\land$
  (conjunção), $\lor$ (disjunção), $\lif$ (condicional)
\item Os sinais de pontuação: (, ) e a vírgula.
\end{enumerate}

\begin{defn}[Fórmula]
\ollabel{defn:formulas} O conjunto~$\Frm[L_0]$ das \emph{!!{formula}s} da
lógica proposicional intuicionista é definido indutivamente da seguinte maneira:
\begin{enumerate}
\item $\lfalse$ é uma !!{formula} atômica.

\item Toda !!{propositional variable}~$\Obj p_i$ é uma
  !!{formula}.

\item Se $!A$ e $!B$ são !!{formula}s, então $(!A \land
  !B)$ é !!a{formula}.

\item Se $!A$ e $!B$ são !!{formula}s, então $(!A \lor !B)$
  é !!a{formula}.

\item Se $!A$ e $!B$ são !!{formula}s, então $(!A \lif !B)$
  é !!a{formula}.

\tagitem{limitClause}{Nada mais é !!a{formula}.}{}
\end{enumerate}
\end{defn}

Além dos conectivos primitivos introduzidos acima, também usamos os
seguintes símbolos \emph{definidos}: $\lnot$ (negação) e $\liff$
(!!{biconditional}). As fórmulas construídas com os operadores definidos
devem ser entendidas da seguinte maneira:
\begin{enumerate}
\item $\lnot !A$ abrevia $!A \lif \lfalse$.

\item $!A \liff !B$ abrevia $(!A \lif !B) \land (!B
  \lif !A)$.
\end{enumerate}

Embora $\lnot$ seja oficialmente tratada como abreviação, às vezes
apresentaremos regras explícitas e cláusulas de definições para~$\lnot$
como se fosse primitiva. Isso serve principalmente para podermos
formular exercícios.

\end{document}
